% =====================================================================
%  preamble.tex
%  Packages, theorem environments, and styling.
%  Custom math shorthands live in macros.tex.
% =====================================================================

% --- Encoding and fonts (pdflatex-friendly) ---------------------------
\usepackage[utf8]{inputenc}
\usepackage[T1]{fontenc}
\usepackage{lmodern}
\usepackage{microtype}
\usepackage{centernot}
% --- Math -------------------------------------------------------------
\usepackage{amsmath, amssymb, amsthm, mathtools}
\usepackage{mathrsfs}
% --- Graphics, tables, lists -----------------------------------------
\usepackage{graphicx}
\graphicspath{{figures/}}
\usepackage{tikz}
\usepackage{tikz-3dplot}
\usetikzlibrary{arrows.meta,calc,decorations.markings,decorations.pathreplacing}
\usepackage{booktabs}
\usepackage{enumitem}
\usepackage{float}
\usepackage{tabularx}
\usepackage{pgfplots}
\usepgfplotslibrary{colormaps}
% Use high compat level for best features
\pgfplotsset{compat=1.18}
% Define a style for 3D axis handling
\pgfplotsset{
    standard 3d/.style={
        axis lines=center,
        width=8cm, height=8cm,
        xmin=-1.5, xmax=1.5,
        ymin=-1.5, ymax=1.5,
        zmin=0, zmax=2.2,
        ticks=none,
        % Perspective angle
        view={115}{20},
        % Style the labels slightly out
        xlabel style={at={(ticklabel* cs:1)},anchor=north west},
        ylabel style={at={(ticklabel* cs:1)},anchor=south west},
        zlabel style={at={(ticklabel* cs:1)},anchor=south},
    }
}
% --- Page geometry ---------------------------------------------------

% ── Page geometry ───────────────────────────────────────────
\usepackage[
  top=.8in,
  bottom=.8in,
  inner=.8in,
  outer=.8in,
  marginparwidth=0.8in
]{geometry}

% ── Headers and footers ─────────────────────────────────────
\usepackage{fancyhdr}
\pagestyle{fancy}
\fancyhf{}
\fancyhead[LE]{\small\itshape\leftmark}    % chapter title on even pages
\fancyhead[RO]{\small\itshape\rightmark}   % section title on odd pages
\fancyfoot[C]{\thepage}
\renewcommand{\headrulewidth}{0.4pt}

% Redefine plain style (used on chapter-opening pages)
\fancypagestyle{plain}{%
  \fancyhf{}%
  \fancyfoot[C]{\thepage}%
  \renewcommand{\headrulewidth}{0pt}%
}

% ── Part and chapter title formatting (optional) ────────────
% Uncomment and customise if you want non-default styling:
% \usepackage{titlesec}
% \titleformat{\chapter}[display]
%   {\normalfont\huge\bfseries}
%   {\chaptertitlename\ \thechapter}{20pt}{\Huge}

% ── Spacing ─────────────────────────────────────────────────
\usepackage{setspace}
% --- Cross-references and links --------------------------------------
% pdfencoding=unicode lets math in section titles bookmark cleanly.
\usepackage[
  colorlinks=true,
  linkcolor=blue!50!black,
  citecolor=blue!50!black,
  urlcolor=blue!50!black,
  pdfencoding=unicode,
  bookmarksnumbered=true
]{hyperref}
\usepackage[capitalize, nameinlink]{cleveref}

% =====================================================================
%  Theorem-like environments
%  Numbered together within each chapter so cross-refs read 7.2.1, etc.
% =====================================================================

\theoremstyle{plain}
\newtheorem{theorem}{Theorem}[chapter]
\newtheorem{proposition}[theorem]{Proposition}
\newtheorem{lemma}[theorem]{Lemma}
\newtheorem{corollary}[theorem]{Corollary}

\theoremstyle{definition}
\newtheorem{definition}[theorem]{Definition}
\newtheorem{example}[theorem]{Example}
\newtheorem{counterexample}[theorem]{Counterexample}

\theoremstyle{remark}
\newtheorem{remark}[theorem]{Remark}
\newtheorem{warning}[theorem]{Warning}

% Unnumbered narrative environments used by this book's pedagogy.
% "recall" = explicit callback to earlier material.
% "hook"   = end-of-section forward pointer (see project instructions).
\theoremstyle{remark}
\newtheorem*{recall}{Recall}
\newtheorem*{hook}{Looking ahead}

% Cleveref names for the custom envs
\crefname{counterexample}{Counterexample}{Counterexamples}
\Crefname{counterexample}{Counterexample}{Counterexamples}

% --- Custom math commands --------------------------------------------
% =====================================================================
%  macros.tex
%  Mathematical shorthands. Keep this short; add as the book grows.
% =====================================================================

% --- Number sets ------------------------------------------------------
\newcommand{\R}{\mathbb{R}}
\newcommand{\Z}{\mathbb{Z}}
\newcommand{\N}{\mathbb{N}}
\newcommand{\Q}{\mathbb{Q}}
\newcommand{\Cplx}{\mathbb{C}}   % \C is sometimes claimed by other packages

% --- Vectors ----------------------------------------------------------
% Default style: bold roman. Switch to \boldsymbol if you prefer italic.
\newcommand{\vect}[1]{\mathbf{#1}}
\newcommand{\zerovec}{\mathbf{0}}

% --- Partial derivatives ---------------------------------------------
\newcommand{\pd}[2]{\dfrac{\partial #1}{\partial #2}}
\newcommand{\pdmix}[3]{\dfrac{\partial^{2} #1}{\partial #2 \, \partial #3}}

% --- Norms and absolute value ----------------------------------------
\DeclarePairedDelimiter{\norm}{\lVert}{\rVert}
\DeclarePairedDelimiter{\abs}{\lvert}{\rvert}

% --- Differential forms ----------------------------------------------
% Use \dform for the upright "d" if you want forms typeset distinctively.
% (Many authors leave it italic; switch the definition to taste.)
\newcommand{\dform}{\mathrm{d}}

% --- Common operators -------------------------------------------------
\DeclareMathOperator{\grad}{grad}
\DeclareMathOperator{\divg}{div}
\DeclareMathOperator{\curl}{curl}
\DeclareMathOperator{\Hess}{Hess}
\DeclareMathOperator{\Image}{Im}
\DeclareMathOperator{\Kernel}{Ker}
\DeclareMathOperator{\spn}{span}
\DeclareMathOperator{\sgn}{sgn}

% --- Inner products and pairings -------------------------------------
\newcommand{\inner}[2]{\langle #1, #2 \rangle}

\DeclareMathOperator{\Arg}{Arg}